\documentclass[11pt,letterpaper]{article}
\usepackage[margin=0.9in,headheight=15pt]{geometry}
\usepackage{fontspec}
\setmainfont{Latin Modern Roman}
\setmonofont{Latin Modern Mono}[Scale=0.88]
\newfontfamily\chinesefont{Noto Serif CJK SC}
\usepackage{amsmath,amssymb,amsthm,booktabs,longtable,array}
\usepackage[protrusion=false]{microtype}
\usepackage{enumitem,fancyhdr,needspace,xurl}
\usepackage[hidelinks,unicode]{hyperref}
\hypersetup{pdftitle={A Lean verified finite completion of the odd cycle assertion in Erdos Problem 883},pdfauthor={Yicheng Pan},pdfsubject={Formal source aligned research manuscript, 5 October 2026}}
\pagestyle{fancy}\fancyhf{}
\fancyhead[L]{\small Erdős 883\quad The first question}
\fancyhead[R]{\small Formal source aligned edition}
\fancyfoot[C]{\thepage}\renewcommand{\headrulewidth}{0pt}
\setlength{\parindent}{0pt}\setlength{\parskip}{5pt plus 1pt minus 1pt}
\setlength{\emergencystretch}{3em}\allowdisplaybreaks[1]
\setlist{nosep,leftmargin=1.5em}
\numberwithin{equation}{section}
\newtheorem{theorem}{Theorem}
\newtheorem{lemma}[theorem]{Lemma}
\newcommand{\pos}[1]{\max(0,#1)}
\newcommand{\code}[1]{\texttt{#1}}
\begin{document}
\thispagestyle{plain}
\begin{center}
{\LARGE\bfseries A Lean verified finite completion of the odd cycle assertion in Erdős Problem 883\par}
\vspace{10pt}
{\large Yicheng Pan ({\chinesefont 潘奕成})\par}
\vspace{5pt}
{5 October 2026\par}
{\small Research manuscript for author and expert review\par}
\end{center}
\begin{abstract}
For every natural number $n$, a subset of $\{1,\ldots,n\}$ with more than
$\lfloor n/2\rfloor+\lfloor n/3\rfloor-\lfloor n/6\rfloor$ elements has a coprime graph containing every odd cycle of length $2k+1$, $1\le k\le\lfloor n/6\rfloor$.
We describe a Lean 4 proof of this assertion, the first question of Erdős Problem 883. The argument adapts Della Pietra's surplus smoothing, totient profiles, signature ordering and ordered-Hall architecture. An exact all-subsets interval criterion, checked arithmetic certificates and an analytic theorem for $n\ge200000$ remove the unspecified starting threshold. An elementary odd-triple argument supplies the coprime endpoints. The formal conclusion is checked against the unabbreviated canonical statement with no additional hypotheses. This paper explains the proof and identifies its certificate interfaces; the accompanying frozen Lean source is the authoritative formal artifact. Neither external human review nor a second clean rebuild of every numerical certificate is claimed.
\end{abstract}
\textbf{Keywords:} coprime graph; odd cycles; extremal graph theory; exact certificates; Lean; formal verification.

\section{Statement and contribution}\label{sec:main}
Write $[n]=\{1,\ldots,n\}$, with $[0]=\varnothing$. Let $G(A)$ be the simple graph on $A$ whose distinct vertices $u,v$ are adjacent precisely when $\gcd(u,v)=1$.
\begin{theorem}\label{thm:main}
For every natural number $n$ and every $A\subseteq[n]$, if
\[
 |A|>f(n):=\left\lfloor\frac n2\right\rfloor+
 \left\lfloor\frac n3\right\rfloor-\left\lfloor\frac n6\right\rfloor,
\]
then $G(A)$ contains a simple cycle of every odd length $l$ satisfying
$3\le l\le\lfloor n/3\rfloor+1$.
Equivalently, it contains $C_{2k+1}$ for every $1\le k\le\lfloor n/6\rfloor$.
\end{theorem}
The equivalence of the two length ranges is proved in the formal development. In particular there is no ambient-size hypothesis hidden in the statement. For $n<6$ the requested range is empty.

Erdős and Sárközy proved a sufficiently-large-$n$ odd-cycle result with some positive linear half-length constant and proposed $1/6$ as the possible sharp value~\cite{ES97}. Della Pietra's pinned 27 July 2026 manuscript states the sharp half-length range for sufficiently large $n$~\cite{DP26}. Its surplus smoothing, totient-profile analysis, prime-signature ordering and Hall construction are central antecedents of the present argument. Those organizing ideas are not claimed here as new.

The contribution described here is a fully specified finite completion and formal verification of the all-$n$ assertion: a uniform interval criterion for every permitted subset, an explicit analytic starting point of $200000$, exact certificates for the remaining range, and a kernel-checked assembly into the canonical first-question statement. The endpoint input is also replaced by an elementary pigeonhole proof. These are claims about the delivered proof and its refinements, not an assertion of worldwide priority. No exhaustive priority search or live-community verification is represented by this manuscript.

The separate complete-tripartite question is outside Theorem~\ref{thm:main}. Sárközy's 1999 paper concerns that question~\cite{S99}; a public Lean project directed at that second question is likewise distinct~\cite{SecondLean}. Neither is a result of this work.

\subsection{Status of the formal and written proofs}
The principal declaration is
\begin{quote}\small\ttfamily
Erdos883Verified.erdos883\_firstQuestion :\par
Erdos883Target.FirstQuestion
\end{quote}
in \path{Erdos883VerifiedCoverage.lean}. Its axiom report lists exactly \code{propext}, \code{Classical.choice} and \code{Quot.sound}. A separate check types the same theorem at the full canonical proposition, without the named abbreviation. The formal cycle notion uses the canonical graph-walk cycle definition, not a weaker closed-walk surrogate.

This exposition is revised to follow the final formal interfaces. It supersedes the earlier informal draft's $2000000$ analytic threshold, external endpoint assumption and $71+84$ interval inventory. It does not assert that every proof step in the older manuscript is textually or mathematically identical to the final Lean route. Numerical claims below are exact rational inequalities; their large finite witnesses reside in Lean source rather than in rounded output tables.

\subsection{AI assistance and responsibility}
OpenAI generative AI systems were used substantially in mathematical development, proof and certificate programming, formalization, literature investigation, checking, and exposition. This was not limited to language editing. Automated and AI-assisted checks do not constitute independent human mathematical review. The precise tool and model record, the author's personal contributions, and the extent of the author's independent review must be finalized before submission; no unrecorded human contribution or review is asserted here.

\section{The exact interval criterion}\label{sec:criterion}

\subsection{Notation and the elementary endpoint lemma}
Let
\[
\begin{aligned}
f(n)&=\lfloor n/2\rfloor+\lfloor n/3\rfloor-\lfloor n/6\rfloor,
& m(n)&=\lfloor n/6\rfloor,\\
H(n)&=\lceil n/2\rceil,
& M(n)&=\lfloor n/3\rfloor-\lfloor n/6\rfloor+1.
\end{aligned}
\]
Let $O_U$ be the odd integers in $[U]$, in any fixed total order (one may use decreasing $\phi(v)/v$). Write $P_k$ for its first $k$ members.

The retained odd count has the useful identity
\[
 M(n)=\left\lfloor\frac{n+3}{6}\right\rfloor+1.
\]
Any set of $M(n)$ distinct odd integers in $[n]$, for $n\ge6$, contains a coprime pair. Here is a direct proof. If the set contains $1$, choose any other member; $M(n)\ge2$. Otherwise map every member $a$ to $\lfloor(a+3)/6\rfloor$. The image lies in $\{1,\ldots,\lfloor(n+3)/6\rfloor\}$, so two distinct members lie in one fibre. Each fibre consists of three consecutive odd integers, of the form $\{6k+3,6k+5,6k+7\}$. Their pairwise differences are $2$ or $4$; a common divisor of two odd members must therefore be $1$.

This is the entire endpoint input in the formal proof: \texttt{oddEndpointPrinciple} in \path{Erdos883AKEndpoints.lean}. The historical filename is retained, but no Ahlswede--Khachatrian theorem is assumed. The argument replaces the external endpoint input used in the earlier informal manuscript and in the cited asymptotic architecture.

Choose any ordered finite list of divisor coordinates $c_1,\ldots,c_t$. The $s$-signature of $v$ is $(1_{c_i\mid v})_{1\le i\le s}$. The graph-theoretic criterion does not require these coordinates to be prime. Set $h_0=0$ and $h_s=2^s+1$ for $s\ge1$.
For integers $6\le L\le U$, define:
\[
\begin{aligned}
e_s(k;L,U)
&=\min_{\substack{u,v\in P_k,\ u\ne v\\\text{equal $s$-signature}}}
\#\{w\in[L]:w\text{ even},\ \gcd(w,u)=\gcd(w,v)=1\};\\
r_s(k;L,U)
&=\text{the same minimum with all positive }w\le L.
\end{aligned}
\]
If there is no such pair, set the corresponding minimum to $+\infty$. In Lean, \texttt{ExactIntervalCertificate} expresses the corresponding lower bound pointwise for every eligible pair, so it needs no infinite value or numerical sentinel.

\subsection*{Exact interval criterion}
Suppose for every integer
\[
0\le b\le H(U)-M(U),\qquad 1\le j\le m(U),
\]
there exist an $s\in\{0,\ldots,t\}$ and $k\in\{0,\ldots,H(U)\}$ such that
\begin{equation}
2\max(0,H(U)-k-b)+h_s<j, \tag{1}
\end{equation}
and at least one of
\begin{align}
e_s(k;L,U)&\ge b+j, \tag{2E}\\
r_s(k;L,U)&\ge U-f(U)+m(U)+j. \tag{2R}
\end{align}
Then for every integer $n\in[L,U]$ and every $A\subseteq[n]$ with $|A|>f(n)$, $G(A)$ contains every odd cycle $C_{2\ell+1}$, $1\le\ell\le\lfloor n/6\rfloor$.

\subsection*{Proof}
Fix such $n,A$ and put $b=\#(\text{even integers in }[n]\text{ outside }A)$. There are at least $M(n)+b$ selected odds. Retain the first $M(n)$ of these in the fixed $O_U$ order, calling the resulting set $Y_0$. For every $k$,
\begin{equation}
|Y_0\setminus P_k|\le\max(0,H(U)-k-b). \tag{3}
\end{equation}
Indeed, if a retained odd lies after position $k$, every discarded selected odd lies after $k$ as well; there are at least $b$ discarded odds. Counting inside $O_U$ proves~(3). If no retained odd lies after $k$ the assertion is immediate.

By the elementary endpoint lemma, choose coprime distinct $a,a'$ in $Y_0$. Since $M(n)\ge m(n)+1$, for each $1\le\ell\le m(n)$ choose $\ell-1$ other members of $Y_0$. Sort these middle vertices lexicographically by their full $t$-signatures, then place $a$ first and $a'$ last. Denote the result $o_0=a,\ldots,o_\ell=a'$. This is the construction in \texttt{exists\_divisor\_skeleton}.

This one order has, simultaneously for every $1\le s\le t$, at most $h_s$ gaps joining different $s$-signatures: within the sorted middle, every prefix-signature occupies one block and there are at most $2^s-1$ transitions; adjoining the two prescribed endpoints adds at most two. This argument also covers an empty middle. For $s=0$ there are no transitions. At most twice the number in~(3) of the gaps touch an odd outside $P_k$.

Let $R=A$ minus the skeleton and let
\[
D_i=\{w\in R:\gcd(w,o_{i-1})=\gcd(w,o_i)=1\}.
\]
For each rank $j\le\ell$ take $s,k$ from the criterion. Every gap except at most $2\max(0,H(U)-k-b)+h_s$ has both endpoints in $P_k$ and equal $s$-signature.

If (2E) holds, that gap has at least $b+j$ common even neighbors in $[L]$. At most $b$ are missing from $A$, and no even vertex is in the odd skeleton, so $|D_i|\ge j$.

If (2R) holds, it has at least $U-f(U)+m(U)+j$ common neighbors in $[L]$. Removing vertices outside $A$ and in the skeleton costs at most
\[
(n-|A|)+(\ell+1)\le n-f(n)+\ell\le U-f(U)+m(U),
\]
so again $|D_i|\ge j$. The last inequality follows from monotonicity of $n-f(n)+m(n)$; writing $n=6q+r$, it equals $3q+(0,1,1,1,1,2)_r$.
Also $b\le H(n)-M(n)\le H(U)-M(U)$, since $H(n)-M(n)=2q+(-1,0,0,0,0,1)_r$ is nondecreasing.

Consequently fewer than $j$ lists $D_i$ have size less than $j$. Sort the $\ell$ lists by increasing cardinality. The $j$th list has size at least $j$, so greedy selection supplies distinct representatives $w_i$. They all lie outside the skeleton. Hence
\[
o_0,w_1,o_1,\ldots,w_\ell,o_\ell,o_0
\]
is a simple coprime cycle of length $2\ell+1$, closing via $\gcd(a,a')=1$. Different ranks may use different $s,k$ or (2E)/(2R), because each bounds the same actual resource lists $D_i$. This proves the criterion.

\section{Arithmetic resources}\label{sec:resources}

\subsection*{Notation}
$\rho(v)=\phi(v)/v$. Throughout this part, endpoint pairs $u,v$ are odd integers.
For the initial tail estimates, specialize the coordinates to the first $s$ odd primes $p_1=3,\ldots,p_s$, and write $p=p_{s+1}$. The uniform construction below will instead use a specified coherent family of divisor-coordinate lists.
For $s=0$, set $p_0=2$. The $0$-signature is empty, $p=p_1=3$, and the tail primes $r>p_0$ are exactly the odd primes; all endpoint pairs remain odd.
Let $U$ be an integer upper bound on every skeleton vertex and let $L\le U$.
Let
\[
\begin{aligned}
H&=\lceil U/2\rceil,& q_L&=\lfloor L/2\rfloor,& m&=\lfloor U/6\rfloor,\\
f&=\lfloor U/2\rfloor+\lfloor U/3\rfloor-\lfloor U/6\rfloor,
& M&=f-\lfloor U/2\rfloor+1.
\end{aligned}
\]
For $0<z\le1$ define the strict finite profile count
\[
R_U(z)=\#\{\text{odd }v\le U:\rho(v)<z\}.
\]
Strict inequality ensures that every remaining vertex has profile $\ge z$.
Write $h_0=0$ and $h_s=2^s+1$ for $s\ge1$.

\subsection*{1. Exact adaptive tail factor}
Let $a$ be the greatest nonnegative integer with $p^a\le U$ and put
\[
\eta_s(U)=\left(\frac{p-1}{p}\right)^a.
\]
Every $v\le U$ has at most $a$ distinct prime divisors exceeding $p_s$. Each corresponding totient factor is at least $(p-1)/p$. Therefore
\[
\prod_{\substack{r\mid v,\ r>p_s\\r\text{ prime}}}\left(1-\frac1r\right)\ge\eta_s(U).
\]
All ingredients are integers or rational numbers; logarithms are unnecessary for computing $a$.

If odd integers $u,v\le U$ share their first $s$ prime-divisibility signatures, every prime factor of $v$ absent from $u$ exceeds $p_s$. Hence
\[
\begin{aligned}
\rho(\operatorname{lcm}(u,v))
&=\rho(u)\prod_{\substack{r\mid v,\ r\nmid u\\r\text{ prime}}}\left(1-\frac1r\right)\\
&\ge\rho(u)\eta_s(U).
\end{aligned}
\]
Interchanging $u,v$ proves the slightly stronger bound
\[
\rho(\operatorname{lcm}(u,v))\ge\max(\rho(u),\rho(v))\eta_s(U).
\]
There is also the signature-free inequality
\[
 \rho(\operatorname{lcm}(u,v))\ge\rho(u)\rho(v),
\]
because the Euler factors at common prime divisors appear twice on the right and once on the left. Thus when both profiles are at least $z\ge0$ the useful combined bound is
\[
 B_s(z):=\max\{z^2,z\eta_s(U)\}\le\rho(\operatorname{lcm}(u,v)).
\]
The final finite adaptive certificates use this combined bound.

One may sharpen $\eta_s(U)$ without changing the proof. Let $r$ be maximal such that the product of the first $r$ primes after $p_s$ is $\le U$. No integer $\le U$ has more than $r$ distinct tail prime factors. Among any $r$ such primes the smallest $r$ give the smallest totient-factor product; fewer factors give no smaller a product. Thus their factor product is a valid sharper rational $\eta_s(U)$. The formal \texttt{SharpTailCertificate} packages this argument as a finite set $Q$ and threshold $q$: all odd primes below $q$ belong either to the signature list or to $Q$, every $x\in Q$ satisfies $2\le x\le q$, and
\[
 U<q\prod_{x\in Q}x,\qquad
 \eta_s(U)=\prod_{x\in Q}(1-1/x).
\]
A product-barrier lemma and an Euler-factor comparison prove the resulting tail bound. In the supplied certificates $Q$ consists of the initial primes beyond the signature prefix.

\subsection*{2. Exact common-neighbor discrepancy}
For odd squarefree $D$ with $w=\omega(D)\ge1$ and integer $X\ge0$,
\[
\begin{aligned}
C_D(X)&=\#\{1\le a\le X:\gcd(a,D)=1\}\\
&=\sum_{d\mid D}\mu(d)\lfloor X/d\rfloor.
\end{aligned}
\]
Comparing with $X\phi(D)/D$, only the positive-$\mu$ terms can contribute negatively. There are $2^{w-1}$ such divisors, and $d=1$ has zero fractional part. Thus
\[
C_D(X)\ge X\rho(D)-(2^{w-1}-1).
\]
For $D=1$ the equality $C_1(X)=X$ holds.

Define $W(U)$ as the greatest $w$ such that the product of the first $w$ odd primes is $\le U^2$, and $\Delta(U)=2^{W(U)-1}-1$ (take zero if $W(U)=0$).
For $D=\operatorname{rad}(uv)$, $D\le uv\le U^2$, so the previous inequality gives
\[
\begin{aligned}
\text{common even neighbors in }[L]&\ge q_L\rho(D)-\Delta(U),\\
\text{common whole-pool neighbors in }[L]&\ge L\rho(D)-\Delta(U).
\end{aligned}
\]
This does not require $D$ to divide $L$ and remains valid if an endpoint is $1$.

\subsection*{3. Exact finite criterion using only a univariate profile}
For every required pair $0\le b\le H-M$ and $1\le j\le m$, suppose one can choose a signature prefix $s$ and rational $z$ with
\begin{equation}
2\max(0,R_U(z)-b)+h_s<j, \tag{A}
\end{equation}
and either
\begin{equation}
q_LB_s(z)-\Delta(U)\ge b+j, \tag{$B_E$}
\end{equation}
or
\begin{equation}
LB_s(z)-\Delta(U)\ge U-f+m+j. \tag{$B_R$}
\end{equation}
Then the all-subsets conclusion of the finite interval criterion holds for every $n\in[L,U]$.

\noindent\textit{Proof:} retain the best $M(n)$ selected odds as before. Among the $R_U(z)$ low profile odds, at most $\max(0,R_U(z)-b)$ can remain in that retained set, because at least $b$ selected odds worse than any retained low-profile odd were discarded. Choose a single full coordinate list containing all prefixes used in the certificate, sort the middle of the skeleton by its full signature, and adjoin its coprime endpoints. For any selected prefix $s$, all but the number in (A) of its gaps have equal signatures and both endpoint profiles $\ge z$. The tail and discrepancy bounds give at least $j$ resources for each such gap. The ordered-Hall proof applies unchanged.

Different ranks may select different prefix lengths. A single ordering by the full list has all shorter prefix signatures contiguous in the middle, so there is no consistency problem.

For the finite implementation, define integer lower degrees
\[
 E_s(v)=\lfloor q_LB_s(\rho(v))-\Delta(U)\rfloor_+,
 \qquad
 Q_s(v)=\lfloor LB_s(\rho(v))-\Delta(U)\rfloor_+,
\]
where $\lfloor y\rfloor_+=\max(0,\lfloor y\rfloor)$. Both are nondecreasing functions of the same profile $\rho(v)$. For each $(b,j)$ it suffices to select $s$ with $h_s<j$ and verify either
\[
 \#\{v\text{ odd}\le U:E_s(v)<b+j\}
 \le b+\left\lfloor\frac{j-h_s-1}{2}\right\rfloor,
\]
or
\[
 \#\{v\text{ odd}\le U:Q_s(v)<U-f+m+j\}
 \le b+\left\lfloor\frac{j-h_s-1}{2}\right\rfloor.
\]
One common decreasing-profile order turns each good-degree set into an initial segment. Its complementary low count gives the strict exceptional-gap budget, while the degree threshold gives the corresponding common-neighbor bound. This is \texttt{DegreeIntervalCertificate}; its reduction to the exact criterion is proved by \texttt{exact\_interval\_of\_degree\_certificate} in \path{Erdos883DegreeCriterion.lean}.

\subsection*{4. A fully uniform tail bound with $O(\sqrt U)$ transitions}
For every integer $U>1024$ choose the least $s$ with $U\le4^s$. Thus $s\ge6$ and $4^{s-1}<U$, whence $2^s<2\sqrt U$.

Here the precise coordinate family matters. Let $\mathcal P$ be the checked list of the first $511$ odd primes, ending at $3671$. Define $\mathcal C_s$ to be the first $s$ entries of $\mathcal P$ when $s\le511$; for $s>511$, append $s-511$ consecutive odd integers beginning with $3673$. These appended coordinates need not be prime. The family has length $s$ and is coherent under taking prefixes. For $s\ge511$, it contains every odd prime below $2s+2651$. The source checks the coverage of the finite list by certifying that omitted integers in $[3,3671]$ have proper divisors.

If $u,v$ share their $\mathcal C_s$-signatures, every prime factor of $v$ absent from $u$ is at least a suitable threshold $p$. Let $a$ count these missing prime factors; then $p^a\le v\le U\le4^s$.

For $s\ge511$, take $p=2s+2651$, so $p\ge2s+3\ge1024$. If $10a\ge p$, then
\[
 4^s=2^{2s}<2^{10a}\le p^a,
\]
a contradiction. Therefore $10a<p$, and Bernoulli's inequality gives
\[
 \prod_{\substack{r\mid v,\ r\nmid u\\r\text{ prime}}}\left(1-\frac1r\right)
 \ge(1-1/p)^a\ge1-a/p>\frac9{10}.
\]
No estimate for the growth of primes is needed.

For $9\le s<511$, let $p$ be entry $s+1$ of $\mathcal P$. The $502$ kernel-checked integer inequalities
\[
 4^s<p^{\lceil p/10\rceil}
\]
again imply $10a<p$. For $s=6,7,8$, separate exact product barriers give the same strict Euler-product bound: at scale $6$ there are at most two missing primes, and if $19$ occurs the other is at least $23$; at scale $7$ there are at most two, all at least $23$; at scale $8$ there are at most three, all at least $29$. The inequalities $(18/19)(22/23)>9/10$, $(22/23)^2>9/10$, and $(28/29)^3>9/10$ complete these cases.

Consequently,
\[
 h_s=2^s+1<2\sqrt U+1,\qquad
 \rho(\operatorname{lcm}(u,v))>\frac9{10}\max(\rho(u),\rho(v))
\]
whenever $u,v\le U$ have equal full signatures. The declarations are \texttt{adaptiveTailCoordinates}, \texttt{uniformTailScale}, and \texttt{totientDensity\_lcm\_gt\_nine\_tenths\_uniform} in \path{Erdos883FiniteTailThresholds.lean}, with the elementary large-scale argument in \path{Erdos883UniformTail.lean} and the three small cases in \path{Erdos883SmallTailCertificates.lean}.

A useful uniform discrepancy bound is also elementary. For odd squarefree $D$,
\[
\frac{(2^{\omega(D)})^4}{D}
=\prod_{p\mid D}\frac{16}{p}
\le\frac{16^5}{3\cdot5\cdot7\cdot11\cdot13}<81,
\]
since factors from primes $\ge17$ are $\le1$. Hence $2^{\omega(D)}<3D^{1/4}$.
For $D\le U^2$, the actual discrepancy from the discrepancy bound is strictly less than $(3/2)\sqrt U$. Thus $(B_E)/(B_R)$ may use loss $(3/2)\sqrt U$ and factor $9/10$, without computing $W(U)$.

This proves a universal numerical reduction for all $U>1024$:
\[
2\max(0,R_U(z)-b)+2^s+1<j,
\]
and either
\[
\frac9{10}q_Lz-\frac32\sqrt U\ge b+j,
\]
or
\[
\frac9{10}Lz-\frac32\sqrt U\ge U-f+m+j
\]
are sufficient. The exact integer $2^s+1$ is preferable to its displayed $2\sqrt U+1$ upper bound. Irrational square roots may be rounded upward in the subtracted error if a rational-only certificate is desired.


\section{Finite distribution estimates}\label{sec:cdf}
For $n>0$ and $t>0$ put
\[
 R_n(t)=\#\{v\le n:v\text{ odd},\ \rho(v)<t\},\qquad
 F_n(t)=\frac2n R_n(t).
\]
The normalization is by $n/2$, not by the number of odd integers. This distinction matters when $n$ is odd. The formal development establishes the following three bounds:
\begin{align}
 F_n(t)&<22t^6 &&(n>0,\ t>0),\label{eq:moment}\\
 F_n(t)&<\frac9{10}t-\frac16-\frac{11}{500}
 &&\left(n\ge200000,\ \frac13\le t\le\frac34\right),\label{eq:cdfhigh}\\
 F_n\left(\sqrt{\frac{\beta+9/125}{3}}\right)
 &<\frac\beta3-\frac1{400}
 &&\left(n\ge200000,\ \frac1{10}\le\beta\le\frac{201}{100}\right).
 \label{eq:cdflow}
\end{align}
The sixth-moment statement is \code{odd\_totient\_cdf\_lt\_twenty\_two\_real\_uniform}; the latter two are \code{totientCdf\_refined\_high} and \code{totientCdf\_refined\_low}. The word ``finite'' here means estimates for each actual finite prefix of the integers. Replacing these by a limiting distribution without an explicit finite-prefix error would not justify the proof.

\subsection{The global sixth moment}
For an odd prime $p$ set $W_p=(p/(p-1))^6-1$. Expanding the finite divisor product for $\rho(v)^{-6}$ and bounding the number of multiples of each divisor by $n/d$ gives
\[
 \sum_{\substack{v\le n\\v\text{ odd}}}\rho(v)^{-6}
 \le n\prod_{\substack{p\le n\\p\text{ odd prime}}}\left(1+\frac{W_p}{p}\right)<11n.
\]
The constant is certified as follows. The exact finite product over odd primes at most $200$, multiplied by $200/193$, is less than $11$. Beyond $200$, $W_p\le7/(p-1)$, so the sum of the positive product increments is at most $7/200$. The inequality $\prod(1+x_i)\le1/(1-\sum x_i)$ bounds the tail by $200/193$. All head arithmetic and tail inequalities are proved in the closure rooted at \path{Erdos883EulerBound.lean}. Each term with $\rho(v)<t$ exceeds $t^{-6}$, giving~\eqref{eq:moment}. This argument normalizes a bound by $n$, not an unjustified bound by the odd count $H(n)$.

\subsection{An exact rational envelope}
The refined envelope splits the odd integers by divisibility by $3,5,7,11$. Start with the state $(a,q,e)=(1,1,1/2)$. For each of these four primes $p$, replace every state by the two states
\[
 \left(a,q\frac{p-1}{p},2e\right),\qquad
 \left(a\frac{p-1}{p},\frac qp,e\right).
\]
There are $16$ final states. Here $a$ is the head totient density, $q$ its normalized asymptotic mass, and $e$ its finite counting-error parameter. For each state define
\[
 c=q+\frac{2e}{200000},\qquad
 d=\frac{47q}{2000}+\frac{991e}{100\cdot200000},
\]
and, for $t>0$,
\[
 B_{a,q,e}(t)=
 \begin{cases}
 c,&a\le t,\\
 \min\left(c,\ d\dfrac{a+t}{2(a-t)}\right),&t<a.
 \end{cases}
\]
Set
\[
 B(t)=\min\left(22t^6,\sum_{(a,q,e)}B_{a,q,e}(t)\right).
\]
The formal envelope theorem proves $F_n(t)\le B(t)$ for every $n\ge200000$ and $t>0$. The cap follows from finite residue-class counting. For $t<a$, a statewise logarithmic tail estimate and
\[
 \log(a/t)\ge\frac{2(a-t)}{a+t}
\]
give the rational second term. The finite-prefix logarithmic error is controlled uniformly by
\[
 \frac{\log(n/10)}n\le\frac{991/100}{200000},
 \qquad n\ge200000.
\]
More explicitly, if $K_S$ counts a state and $K_{S,p}$ counts its members divisible by an additional prime $p>11$, residue-class inclusion--exclusion gives
\[
 |K_S-nq/2|\le e,\qquad |K_{S,p}-nq/(2p)|\le e.
\]
With $\lambda_p=\log(p/(p-1))$, summing the tail divisibility indicators bounds the normalized state log moment by
\[
 \frac{47q}{2000}+\frac{e\log(n/10)}{n}.
\]
The prime sum $\sum_{p>11}\lambda_p/p<47/2000$ follows from the rational majorant $1/p^2+1/(2p^2(p-1))$: the exact head through $2000$ is at most $22968311/10^9$, and the tail is at most $1/2000$. Dividing the log-moment bound by $\log(a/t)$ gives the state tail estimate. These arithmetic bounds and their finite counting interfaces are proved in the dependency closure, rather than imported as numerical assumptions. The global sixth-moment bound gives the other arm of the minimum. The formal bridge from the actual state counts to this envelope is in \path{Erdos883RefinedCdfEnvelope.lean}; the rational state definition is in \path{Erdos883RefinedCdfDefinitions.lean}.

\subsection{What the finite CDF certificates check}
The high-range certificate has $1000$ cells. With $t_i=1/3+i/2400$, it verifies, for $0\le i<1000$,
\[
 \frac9{10}t_i-\frac16-B(t_{i+1})>\frac{11}{500}.
\]
Monotonicity and endpoint handling then give~\eqref{eq:cdfhigh} on the entire closed interval, not just at sampled points.

The low-range certificate has $1910$ cells. Put $\beta_i=1/10+i/1000$. For each $0\le i<1910$ the source supplies an integer $z_i$ and verifies, for $u_i=z_i/10^9$,
\[
 u_i\ge0,\quad
 \frac{\beta_{i+1}+9/125}{3}\le u_i^2,\quad
 \frac{\beta_i}3-B(u_i)>\frac1{400}.
\]
Thus the square-root upper witnesses are checked by exact squaring. Their origin need not be trusted. Monotonicity extends each cell inequality to the full interval and yields~\eqref{eq:cdflow}, including the endpoint $201/100$. In total these refined CDF checks comprise $2910$ rational cells. They are proved in the \code{Erdos883RefinedCdfCells*} modules and interpolated in \code{Erdos883RefinedCdfInterpolation*}. No floating-point pass criterion occurs in this interface.

These finite inequalities form a numerical lemma with machine-readable proof, not a substitute for the reduction to them. The preceding envelope explains what is counted and why the cells control every real threshold needed below. An expert may review the conceptual reduction separately from replaying the exact certificates.

\section{The analytic range starting at 200000}\label{sec:analytic}
Fix $n\ge200000$ and put
\[
 m=\lfloor n/6\rfloor,\quad \delta=\frac32\sqrt n,\quad
 \varepsilon=\frac{\delta+2}{m},\quad
 \beta=\frac bm,\quad x=\frac jm.
\]
Here $0\le b\le H(n)-M(n)$ and $1\le j\le m$ are the missing-even and Hall-rank parameters. Let $s$ be the least integer with $n\le4^s$ and put $h=2^s+1$. The final accounting lemmas give
\begin{align}
 \varepsilon&<\frac{101}{5000},&
 \frac hm&<\frac{269}{10000},&
 \frac{n}{2m}&<\frac{37501}{12500},\label{eq:errors}\\
 n-f(n)+m&\le3m+2,&
 b&\le2m+1.\label{eq:rounding}
\end{align}
They retain all rounding errors. For example $6m\le n\le6m+5$, and the square-root bounds can be reduced to quadratic inequalities using $\sqrt n\ge447$.

Write $r=R_n(t)/m$ at a selected threshold. Equations~\eqref{eq:moment} and~\eqref{eq:errors} imply
\[
 r<\frac{66003}{1000}t^6.\tag{M}
\]
Whenever $F_n(t)\le1$, the normalization also gives
\[
 r<3F_n(t)+\frac1{12500}.\tag{N}
\]
The resource estimates from Section~\ref{sec:resources} take particularly simple forms. For a same-signature pair with both profiles at least $t$, the even and whole-pool common-neighbor counts are respectively strictly greater than
\[
 \frac{27}{10}mt-\delta,\qquad
 \frac{27}{5}mt-\delta.\tag{S}
\]
Without any signature condition they are strictly greater than
\[
 3mt^2-\delta,\qquad6mt^2-\delta.\tag{U}
\]
The adaptive coordinate list used for (S) is a coherent family of divisibility signatures as described in Section~\ref{sec:resources}. Different rank choices therefore still act on one common skeleton order.

\subsection{Large ranks}
Suppose $x\ge1/20$ and choose
\[
 t=\min\left(\frac{\beta+x+\varepsilon}{27/10},
              \frac{3+x+\varepsilon}{27/5}\right).
\]
Then $0<t<3/4$. One of the two arms in (S) supplies the corresponding resource target: the first arm gives more than $b+j$ even neighbors, while the second gives more than $n-f(n)+m+j$ whole-pool neighbors. The extra $2$ in $m\varepsilon=\delta+2$ accommodates~\eqref{eq:rounding}.

It remains to prove the normalized rank budget
\[
 2\pos{r-\beta}+h/m<x.\label{eq:budget}\tag{B}
\]
If $t\ge10/27$, use~\eqref{eq:cdfhigh} and (N), with $t\le(\beta+x+\varepsilon)/(27/10)$, to obtain
\[
 r<\beta+x+\varepsilon-\frac{7074}{12500}.
\]
If $r\le\beta$, (B) follows from $h/m<0.0269<0.05\le x$. Otherwise,
\[
 2(r-\beta)+h/m
 <2x+2\varepsilon-\frac{14148}{12500}+\frac{269}{10000}<x,
\]
since $x\le1$ and $\varepsilon<101/5000$.

If $t<10/27$, the first arm of the minimum must be active. Let $y=\beta+x+\varepsilon$. Then $1/20\le y<1$ and (M) yields
\[
 r<\frac{171}{1000}y^6.
\]
The elementary inequality
\[
 \frac{171}{500}y^6-y+\frac{471}{10000}<0
 \qquad(1/20\le y\le1)
\]
proves (B) using the same error bounds. To check this polynomial without differentiation, split at $y=1/10$: below that point bound $y^6\le10^{-6}$; above it use $y^6\le y$. This is the formal low-profile large-rank branch.

\subsection{Small ranks with many missing evens}
Now let $x<1/20$ and $\beta\ge1/10$. Put $\beta_0=\min(\beta,2)$ and choose no signature coordinates, so the transition cost is zero. Take
\[
 t=\sqrt{\frac{\beta_0+9/125}{3}}.
\]
Equations~\eqref{eq:cdflow} and (N) imply
\[
 r<\beta_0-\frac{371}{50000}<\beta,
\]
where the last strict inequality is understood via $\beta_0\le\beta$. Thus the exceptional-gap budget vanishes. Clipping loses at most one vertex because
\[
 0\le m(\beta-\beta_0)\le1.
\]
The even resource bound (U) is sufficient: its main term is $m\beta_0+(9/125)m$, while $j\le m/20$, $\delta+2<(101/5000)m$, and the clipping loss is at most one. The remaining positive margin supplies $b+j$ neighbors. This branch is the reason the finite-prefix CDF bound is needed beyond the small-$\beta$ region.

\subsection{Small ranks with few missing evens}
Finally suppose $x<1/20$ and $\beta<1/10$. There are two subcases.

If $b+j\le30\sqrt n$, no positive profile threshold is needed. The improved uniform neighbor lemma gives more than $30\sqrt n$ common even neighbors for every pair of positive odd integers at most $n$. Its proof uses
\[
 8000\,\omega(\operatorname{lcm}(u,v))\le n
\]
and a finite-product inequality of Wallis type,
\[
 3\le(4\omega(D)+3)\rho(D)^2\qquad(D\text{ positive and odd}).
\]
For completeness, the support bound is immediate if $\omega(D)\le25$. Otherwise use $3^{\omega(D)}\le D\le n^2$ and $(8000w)^2\le3^w$ for $w\ge25$, proved by induction. The density inequality gives $n\rho(D)^2\ge600000/103$. Since $\lfloor n/2\rfloor\ge49n/100$, the square of the even-neighbor main term is at least $(144060/103)n>(63/2)^2n$. Thus that main term exceeds $(63/2)\sqrt n$. Subtracting the discrepancy $\delta=(3/2)\sqrt n$ leaves the required strict $30\sqrt n$ bound. The source is \path{Erdos883ImprovedNeighbors.lean}.

If $b+j>30\sqrt n$, take no signature coordinates and put
\[
 t=\sqrt{\frac7{20}(\beta+x)}.
\]
Now (M) gives
\[
 r<\frac{71}{25}(\beta+x)^3.
\]
Because $\beta+x<3/20$ and $x\le1/20$, this implies $2\pos{r-\beta}<x$. Indeed, if $x\le\beta$, then $(142/25)(\beta+x)^2<1$ gives $r<\beta$. If $\beta<x$, use $\beta+x<2x$ and $(1136/25)x^2<1$ to obtain $2r<x$.

For the resource target, $\delta<(b+j)/20$, whereas
\[
 3mt^2=\frac{21}{20}(b+j).
\]
Thus (U) supplies more than $b+j$ common even neighbors. This completes all rank and missing-even cases.

\subsection{Closing the analytic argument}
For each $(b,j)$ the selected threshold and coordinate prefix satisfy the exact profile criterion. The real-profile ordering lemma constructs a single decreasing profile order that works for every real threshold, so choices of $t$ do not require separate skeletons. The analytic certificate is \path{Improved.largeN_exact_interval_certificate}. The cycle conclusion is \path{Improved.firstQuestion_largeN}. Both appear in \path{Erdos883ImprovedLargeNCover.lean}. They hold for every $n\ge200000$ with no remaining numerical or distribution hypothesis.

\section{Exact finite completion}\label{sec:finite}
The finite computation is a collection of proofs of sufficient arithmetic conditions. It does not sample subsets $A$ and does not assume that behavior observed below a cutoff extrapolates above it. The exact interval criterion in Section~\ref{sec:criterion} is quantified over all permitted subsets and all target ranks.

\subsection{Certificate semantics}
For odd endpoints with profiles at least $t$, both the unstructured product bound and a same-signature tail bound are available. Their combination is
\[
 \rho(\operatorname{lcm}(u,v))\ge
 g_s(t):=\max\{t^2,\eta_s(U)t\}.
\]
At interval endpoints $L\le U$, exact lower-degree functions can be formed from
\[
 \left\lfloor\lfloor L/2\rfloor\,g_s(\rho(v))-\Delta(U)\right\rfloor_+,
 \qquad
 \left\lfloor L\,g_s(\rho(v))-\Delta(U)\right\rfloor_+.
\]
Here $\lfloor y\rfloor_+=\max(0,\lfloor y\rfloor)$ is the natural-valued floor used by the formal degree function. Counting how many profile rows fail a degree threshold bounds the number of exceptional skeleton vertices. With $a=\lfloor(j-h_s-1)/2\rfloor$, the even resource alternative is sufficient if at most $b+a$ rows fail the target $b+j$; the whole resource alternative uses target $U-f(U)+m(U)+j$. Only prefixes with $h_s<j$ can be useful. The strict rank inequality and the distinction between ``less than a degree target'' and ``at most a degree target'' are retained in the certificate definitions.

The Lean finite route checks more than a list of numerical PASS outputs. The certificate interfaces establish that profile metadata represents the actual totient density; that the odd-vertex rows form the full universe without duplication; that prime support and sharp tail witnesses are valid; that cached count-tree data has the required semantics; and that the rank witnesses satisfy the sufficient criterion. The span certificates compress intervals of profile-row ranks, with the span-checking soundness theorem proving the required coverage throughout each interval. These profile ranks are distinct from the Hall-rank parameter $j$. Consequently external generators and proposed square-root, factor or rank witnesses are outside the trusted proof boundary: their generated claims must be proved by the checked interfaces.

Representative interfaces include \path{AdaptiveProfileRowsValid}, \path{SharpTailCertificate}, and \path{AdaptiveTreeRepresents}, together with rank and span soundness theorems. The files \path{Erdos883AdaptiveCertificateAssembly1226.lean} and \path{Erdos883AdaptiveSpanAssembly199999.lean} show the complete assembly at the first and last adaptive endpoints. At the latter endpoint the formal proof uses seven prime coordinates $3,5,7,11,13,17,19$ and a verified joint-prime-support bound of $9$. These finite coordinate lists should not be confused with the analytic coordinate construction for arbitrarily large $n$.

\subsection{Coverage without gaps}
There are $152$ nonvacuous finite certificate intervals: $64$ direct and $88$ adaptive. The frozen finite prefix proves the statement for $0\le n\le1211$. In the final root, a theorem through $680$ is joined with direct certificate intervals
\[
 [681,749],\ [750,825],\ [826,908],\ [909,999],\ [1000,1100],\ [1101,1211].
\]
The range $1212\le n\le199999$ is covered by $88$ adaptive intervals. Appendix~\ref{app:intervals} lists their upper endpoints; each lower endpoint is one more than its predecessor, starting at $1212$. The final theorem \code{firstQuestion\_verifiedPrefix} proves the canonical assertion on $[0,199999]$ by repeated interval composition.

The analytic range begins at $200000$. Thus the completed partition is
\[
 [0,1211]\ \cup\ [1212,199999]\ \cup\ [200000,\infty).
\]
There is neither an omitted integer nor an unresolved tail. The final declaration applies the proved initial-interval-to-global theorem, whose other branch is the analytic result. This proves Theorem~\ref{thm:main}.

\section{Formal statement and verification boundary}\label{sec:formal}
The canonical target is the positive assertion in part (i) of the Formal Conjectures statement at revision \path{89294ea02bd7cd678d59984add52cb4baef3dbf4}~\cite{FC883}. The explicit checked type is
\begin{quote}\small\ttfamily
\ensuremath{\forall} (n : \ensuremath{\mathbb N}) (A : Finset \ensuremath{\mathbb N}),\par
\quad A \ensuremath{\subseteq} Finset.Icc 1 n \ensuremath{\to}\par
\quad n / 2 + n / 3 - n / 6 < A.card \ensuremath{\to}\par
\quad \ensuremath{\forall} l : \ensuremath{\mathbb N}, Odd l \ensuremath{\to} 3 \ensuremath{\le} l \ensuremath{\to} l \ensuremath{\le} n / 3 + 1 \ensuremath{\to}\par
\quad l \ensuremath{\in} ((SimpleGraph.fromRel Nat.Coprime).induce\par
\qquad (A : Set \ensuremath{\mathbb N})).oddCycleLengths
\end{quote}
All divisions here are natural-number division. The induced graph enforces the vertex set $A$; the canonical cycle definition requires a genuine cycle. The statement check is separate from identifying the name of the theorem. It establishes that there is no additional endpoint, numerical-certificate, density-distribution or lower-bound hypothesis in the theorem's interface.

\subsection{What was checked}
The frozen local import closure has $6831$ source modules, totaling $144287247$ source bytes. The toolchain is Lean $4.33.1$ with mathlib revision \path{0df444a360eaa60ab8c11dca51a86af692955474}. The recorded audit combines modular source compilation, targeted recompilation of concrete source/object correspondence concerns, a fresh final-root compilation, explicit canonical-type checking, and source, object and configuration hash comparisons against the frozen snapshot. The final source scan found no admitted proof, custom axiom, native-evaluation proof shortcut or kernel-checking bypass in the local closure.

The reported assumptions are exactly the three standard Lean axioms: \code{propext}, \code{Classical.choice} and \code{Quot.sound}. This is an ordinary Lean-kernel result relative to the pinned toolchain and dependencies; it is not a separately implemented kernel verification or a bootstrap of Lean and mathlib from source.

\subsection{What is not claimed}
The recorded audit is not a second clean rebuild of all $6831$ numerical and supporting modules. Some earlier successful compilations were observed during construction without durable per-module pre/post hash attestations. Empty logs, completion markers and timestamps alone are not represented as frozen-input proof of a successful rebuild. The audit records how concrete correspondence concerns were resolved. A full clean replay is a useful further reproducibility check and remains distinct from the checks already reported.

Likewise, neither this manuscript nor the audit records independent human mathematical review. Kernel checking establishes the formal proposition, subject to the stated trusted dependencies. Human review still matters for the correspondence between the intended problem and its statement, intellectual attribution, intelligibility, author responsibility and publication standards. The accompanying reproduction guide separates these tasks.

\subsection{Artifact identification and availability}
The authoritative frozen formal artifact is
\path{Erdos883_FirstQuestion_Lean_20261005.zip}, version 5, $20359760$ bytes, with SHA-256
\begin{quote}\small\ttfamily
8f017032cb93f2f0891f6f8200390c161aca\par
bf112dac72567462e642b91ede95
\end{quote}
(the two displayed lines are one hash). It includes the exact source closure, pinned configuration, build order, verification script, canonical harness, provenance manifests, reports and logs. The archive is supplied with this research package; this manuscript does not represent it as already publicly deposited or submitted. A public immutable repository or archival identifier should be added only after an authorized release has occurred.

The earlier journal-preparation manuscript and its historical informal certificate bundle remain prior versions. Their different thresholds and interval lists are not evidence for the new formal inventory. The new source package contains an explicit revision note so that an expert can distinguish the exposition's history from the final formal result.

\section{Conclusion}
The first question of Erdős Problem 883 is proved by the named Lean theorem with its exact canonical all-$n$ type. The proof combines an elementary endpoint lemma, Della Pietra's central smoothing and ordering architecture, exact resource criteria, a finite certificate cover and an analytic range beginning at $200000$. The mathematical and formal claims can be evaluated without assuming novelty, human peer review or a clean replay that has not been performed. The next scholarly steps are author approval of the attribution and contribution record, expert review, and a stable authorized public release of the frozen proof and its explanatory material.

\clearpage
\appendix
\section{The adaptive interval endpoints}\label{app:intervals}
Each row lists consecutive upper endpoints. Read left to right, then down. The first interval starts at $1212$; each subsequent interval starts one more than the preceding endpoint. The list is extracted from the $88$ adaptive pieces of \path{Erdos883VerifiedCoverage.lean}. The source package also includes a machine-readable CSV with both endpoints.
\begin{center}\begin{tabular}{rrrr}\toprule
1226 & 1241 & 1262 & 1286 \\
1310 & 1333 & 1364 & 1399 \\
1435 & 1471 & 1508 & 1546 \\
1585 & 1625 & 1666 & 1708 \\
1751 & 1795 & 1840 & 1887 \\
1935 & 1984 & 2034 & 2085 \\
2138 & 2192 & 2244 & 2300 \\
2358 & 2417 & 2478 & 2540 \\
2604 & 2670 & 2737 & 2874 \\
3018 & 3169 & 3328 & 3495 \\
3670 & 3854 & 4047 & 4250 \\
4463 & 4687 & 4922 & 5169 \\
5428 & 5700 & 5986 & 6585 \\
7244 & 7969 & 8767 & 9644 \\
10609 & 11671 & 12839 & 14124 \\
15537 & 17091 & 18801 & 20682 \\
22751 & 25027 & 27530 & 30284 \\
33313 & 36645 & 40310 & 44342 \\
48777 & 53655 & 59021 & 64924 \\
71417 & 78559 & 86416 & 95058 \\
104564 & 115021 & 126524 & 139177 \\
153095 & 168405 & 185246 & 199999 \\
\bottomrule\end{tabular}\end{center}


\begingroup\small
\begin{thebibliography}{9}
\bibitem{ES97}
P. Erdős and G. N. Sárközy, \emph{On cycles in the coprime graph of integers}, Electronic Journal of Combinatorics 4(2) (1997), R8. DOI: \href{https://doi.org/10.37236/1323}{10.37236/1323}. Primary text: \url{https://emis.de/ft/6615}.
\bibitem{DP26}
D. Della Pietra, \emph{Odd cycles at the sharp threshold in the coprime graph}, unrefereed public manuscript, 27 July 2026 publication revision. Pinned source:
\url{https://github.com/donalddellapietra/erdos-883-proof/blob/1f26276d7850fc9afe68e9d7edecf24aa7cc597d/paper/main.tex}.
Its main theorem assumes $n\ge n_0$ for an unspecified absolute $n_0$. The surplus smoothing, totient-profile, signature-ordering and Hall architecture is credited to this work.
\bibitem{S99}
G. N. Sárközy, \emph{Complete tripartite subgraphs in the coprime graph of integers}, Discrete Mathematics 202 (1999), 227--238. DOI: \href{https://doi.org/10.1016/S0012-365X(98)00359-8}{10.1016/S0012-365X(98)00359-8}. \url{https://web.cs.wpi.edu/~gsarkozy/Cikkek/12.pdf}.
\bibitem{FC883}
Formal Conjectures Authors, \emph{Erdős Problem 883}, statement source at revision \path{89294ea02bd7cd678d59984add52cb4baef3dbf4}.
\url{https://github.com/google-deepmind/formal-conjectures/blob/89294ea02bd7cd678d59984add52cb4baef3dbf4/FormalConjectures/ErdosProblems/883.lean}.
This is the canonical statement reference, not a proof source for the result.
\bibitem{SecondLean}
Public repository \emph{jsp-000734-erdos883-lean}, separate second-question formalization. \url{https://github.com/baobingzhang/jsp-000734-erdos883-lean}. Mentioned to distinguish scope; no claim about its authorship or result is made on behalf of the present author.
\end{thebibliography}
\endgroup

\end{document}
